% Part: first-order-logic
% Chapter: syntax-and-semantics
% Section: formation-sequences

\documentclass[../../../include/open-logic-section]{subfiles}

\begin{document}

\olfileid{fol}{syn}{fseq}

\olsection{构造序列}

通过归纳定义来定义公式，并辅以通过归纳法证明公式性质的相应技术，是一种优雅而高效的方法。
然而，考虑一种更自底向上、逐步构造公式的方法也是有用的，这里我们将通过\emph{构造序列}的概念来介绍这种方法。
%
为了展示如何以这种方式引入项和公式而无需依赖其归纳定义，我们首先引入由某个语言~$\Lang L$ 的符号构成的任意字符串这一概念。

\begin{defn}[字符串]
\ollabel{defn:string}
设 $\Lang L$ 是一阶语言。\emph{$\Lang L$-字符串}是由~$\Lang L$ 的符号组成的有限序列。
当语言~$\Lang L$ 在上下文中已明确时，我们通常将 $\Lang L$-字符串简称为\emph{字符串}。
\end{defn}

\begin{ex}
对于任意一阶语言 $\Lang L$，所有的
$\Lang L$-公式都是 $\Lang L$-字符串，但反之则不然。
例如，\[)(\Obj v_0\lif\lexists\] 是一个
$\Lang L$-字符串，但不是一个 $\Lang L$-公式。
\end{ex}

\begin{defn}[项的构造序列]
\ollabel{defn:fseq-trm}
项 $t$ 的\emph{构造序列}是一个 $\Lang L$-字符串的有限序列 $\tuple{t_0,\dotsc,t_n}$，
满足 $t \ident t_n$，且对所有
$i \leq n$，要么 $t_i$ 是变元或常项符号，要么 $\Lang
L$ 包含一个 $k$ 元函数符号~$f$ 且存在
$m_0,\dotsc,m_k < i$ 使得 $t_i \ident f(t_{m_0},\dotsc,t_{m_k})$。
在需要区分时，我们将项的构造序列称为\emph{项构造序列}。
\end{defn}

\begin{ex}
序列
\[
    \tuple{\Obj c_0, \Obj v_0, \Atom{\Obj f^2_0}{\Obj c_0, \Obj v_0}, \Atom{\Obj f^1_0}{\Atom{\Obj f^2_0}{\Obj c_0, \Obj v_0}}}
\]
是项 $\Atom{\Obj f^1_0}{\Atom{\Obj
f^2_0}{\Obj c_0, \Obj v_0}}$ 的一个构造序列，
\[
    \tuple{\Obj v_0, \Obj c_0, \Atom{\Obj f^2_0}{\Obj c_0, \Obj v_0}, \Atom{\Obj f^1_0}{\Atom{\Obj f^2_0}{\Obj c_0, \Obj v_0}}}.
\]
也是。
\end{ex}

\begin{defn}[公式的构造序列]
\ollabel{defn:fseq-frm}
公式~$!A$ 的\emph{构造序列}是一个 $\Lang L$-字符串的有限序列 $\tuple{!A_0,\dotsc,!A_n}$，
满足 $!A \ident !A_n$，且
对所有 $i \leq n$，要么 $!A_i$ 是原子公式，要么存在
$j,k < i$ 和变元~$x$ 使得下列条件之一成立：
\begin{enumerate}
    \tagitem{prvNot}{$!A_i \ident \lnot !A_j$.}{}%
    \tagitem{prvAnd}{$!A_i \ident (!A_j \land !A_k)$.}{}%
    \tagitem{prvOr}{$!A_i \ident (!A_j \lor !A_k)$.}{}%
    \tagitem{prvIf}{$!A_i \ident (!A_j \lif !A_k)$.}{}%
    \tagitem{prvIff}{$!A_i \ident (!A_j \liff !A_k)$.}{}%
    \tagitem{prvAll}{$!A_i \ident \lforall[x][!A_j]$.}{}%
    \tagitem{prvEx}{$!A_i \ident \lexists[x][!A_j]$.}{}%
\end{enumerate}
在需要区分时，我们将公式的构造序列称为\emph{公式构造序列}。
\end{defn}

\begin{ex}
\[
\tuple{
    \Atom{\Obj A^1_0}{\Obj v_0},
    \Atom{\Obj A^1_1}{\Obj c_1},
    (\Atom{\Obj A^1_1}{\Obj c_1} \land \Atom{\Obj A^1_0}{\Obj v_0}),
    \lexists[\Obj v_0][(\Atom{\Obj A^1_1}{\Obj c_1} \land \Atom{\Obj A^1_0}{\Obj v_0})]
}
\]
是 $\lexists[\Obj v_0][(\Atom{\Obj A^1_1}{\Obj
c_1} \land \Atom{\Obj A^1_0}{\Obj v_0})]$ 的一个构造序列，
\begin{multline*}
\tuple{
    \Atom{\Obj A^1_0}{\Obj v_0},
    \Atom{\Obj A^1_1}{\Obj c_1},
    (\Atom{\Obj A^1_1}{\Obj c_1} \land \Atom{\Obj A^1_0}{\Obj v_0}),
    \Atom{\Obj A^1_1}{\Obj c_1},\\
    \lforall[\Obj v_1][\Atom{\Obj A^1_0}{\Obj v_0}],
    \lexists[\Obj v_0][(\Atom{\Obj A^1_1}{\Obj c_1} \land \Atom{\Obj A^1_0}{\Obj v_0})]
}.
\end{multline*}
也是。
%
从第二个例子可以看出，构造序列
可能包含“废料”：冗余的或不参与构造的公式。
\end{ex}

\begin{prop}\ollabel{prop:formed}
$\Frm[L]$ 中的每个公式~$!A$ 都有一个构造序列。
\end{prop}

\begin{proof}
假设 $!A$ 是原子公式。那么序列 $\tuple{!A}$ 就是~$!A$ 的一个
构造序列。
%
现在假设 $!B$ 和~$!C$ 分别具有构造序列
$\tuple{!B_0,\dotsc,!B_n}$ 和 $\tuple{!C_0,\dotsc,!C_m}$。
%
\begin{enumerate}
  \tagitem{prvNot}{若 $!A \ident \lnot !B$，
      则 $\tuple{!B_0,\dotsc,!B_n,\lnot !B_n}$
      是~$!A$ 的一个构造序列。}{}
  \tagitem{prvAnd}{若 $!A \ident (!B \land !C)$，
      则 $\tuple{!B_0,\dotsc,!B_n,!C_0,\dotsc,!C_m,(!B_n \land !C_m)}$
      是~$!A$ 的一个构造序列。}{}
  \tagitem{prvOr}{若 $!A \ident (!B \lor !C)$，
      则 $\tuple{!B_0,\dotsc,!B_n,!C_0,\dotsc,!C_m,(!B_n \lor !C_m)}$
      是~$!A$ 的一个构造序列。}{}
  \tagitem{prvIf}{若 $!A \ident (!B \lif !C)$，
      则 $\tuple{!B_0,\dotsc,!B_n,!C_0,\dotsc,!C_m,(!B_n \lif !C_m)}$
      是~$!A$ 的一个构造序列。}{}
  \tagitem{prvIff}{若 $!A \ident (!B \liff !C)$，
      则 $\tuple{!B_0,\dotsc,!B_n,!C_0,\dotsc,!C_m,(!B_n \liff !C_m)}$
      是~$!A$ 的一个构造序列。}{}
  \tagitem{prvAll}{若 $!A \ident \lforall[x][!B]$，
      则 $\tuple{!B_0,\dotsc,!B_n,\lforall[x][!B_n]}$
      是~$!A$ 的一个构造序列。}{}
  \tagitem{prvEx}{若 $!A \ident \lexists[x][!B]$，
      则 $\tuple{!B_0,\dotsc,!B_n,\lexists[x][!B_n]}$
      是~$!A$ 的一个构造序列。}{}
\end{enumerate}
根据公式归纳原理，
每个公式都有一个构造序列。
\end{proof}

我们也可以证明其逆命题。这一点很重要，因为它表明我们定义公式的两种方式是等价的：它们给出相同的结果。这也意味着我们可以通过对构造序列的长度施用普通归纳法来证明关于公式的定理。

\begin{lem}
\ollabel{lem:fseq-init}
假设 $\tuple{!A_0,\dotsc,!A_n}$ 是~$!A_n$ 的一个构造序列，
且 $k \leq n$。那么 $\tuple{!A_0,\dotsc,!A_k}$
是~$!A_k$ 的一个构造序列。
\end{lem}

\begin{proof}
练习。
\end{proof}

\begin{prob}
证明 \olref[fol][syn][fseq]{lem:fseq-init}。
\end{prob}

\begin{thm}
\ollabel{thm:fseq-frm-equiv}
$\Frm[L]$ 是所有满足如下条件的 $\Lang L$-字符串~$!A$ 构成的集合：
存在~$!A$ 的一个公式构造序列。
\end{thm}

\begin{proof}
设 $F$ 为语言~$\Lang L$ 中所有具有构造序列的符号串构成的集合。
由 \olref[fol][syn][fseq]{prop:formed} 已知 $\Frm[L] \subseteq F$，故现在只需证明反向包含关系。

假设 $!A$ 具有构造序列 $\tuple{!A_0,\dotsc,!A_n}$。
我们对~$n$ 施用强归纳法来证明 $!A \in \Frm[L]$。
我们的归纳假设是：任何具有长度为 $m < n$ 的构造序列的符号串都属于 $\Frm[L]$。
根据构造序列的定义，$!A \ident !A_n$ 要么是原子公式，要么存在 $j,k < n$ 使得下列情形之一成立：
\begin{enumerate}
    \tagitem{prvNot}{$!A \ident \lnot !A_j$.}{}%
    \tagitem{prvAnd}{$!A \ident (!A_j \land !A_k)$.}{}%
    \tagitem{prvOr}{$!A \ident (!A_j \lor !A_k)$.}{}%
    \tagitem{prvIf}{$!A \ident (!A_j \lif !A_k)$.}{}%
    \tagitem{prvIff}{$!A \ident (!A_j \liff !A_k)$.}{}%
    \tagitem{prvAll}{$!A \ident \lforall[x][!A_j]$.}{}%
    \tagitem{prvEx}{$!A \ident \lexists[x][!A_j]$.}{}%
\end{enumerate}
现在我们分情况讨论。若 $!A$ 是原子公式，则 $!A_n \in \Frm[L_0]$。
否则，假设 $!A \equiv (!A_j \land !A_k)$。
根据 \olref[fol][syn][fseq]{lem:fseq-init}，
$\tuple{!A_0,\dotsc,!A_j}$ 和 $\tuple{!A_0,\dotsc,!A_k}$ 分别是 $!A_j$ 和~$!A_k$ 的构造序列。
由于它们是~$!A$ 的构造序列的真初始子序列，它们的长度均小于~$n$。
因此由归纳假设，$!A_j$ 和~$!A_k$ 属于~$\Frm[L_0]$，再根据公式的定义，$(!A_j \land !A_k)$ 亦然。
其余情形可平行地推理。
\end{proof}

项的构造序列具有与公式的构造序列类似的性质。

\begin{prop}
\ollabel{prop:fseq-trm-equiv}
$\Trm[L]$ 是所有具有项构造序列的 $\Lang L$-字符串 $t$ 构成的集合。
\end{prop}

\begin{proof}
练习。
\end{proof}

\begin{prob}
证明 \olref[fol][syn][fseq]{prop:fseq-trm-equiv}。
提示：使用与 \olref[fol][syn][fseq]{thm:fseq-frm-equiv} 的证明类似的策略。
\end{prob}

构造序列中可能出现两种“废料”：重复的元素，以及与构造该公式或项无关的元素。通过考察最小构造序列，我们可以消除这两类废料。

\begin{defn}[最小构造序列]
\ollabel{defn:minimal-fseq}
公式~$!A$ 的构造序列 $\tuple{!A_0, \ldots, !A_n}$ 是~$!A$ 的\emph{最小构造序列}，如果对~$!A$ 的任何其他构造序列~$s$，$s$ 的长度均大于或等于~$n+1$。

类似地，项~$t$ 的构造序列 $\tuple{t_0, \ldots, t_n}$ 是~$t$ 的\emph{最小构造序列}，如果对~$t$ 的任何其他构造序列~$s$，$s$ 的长度均大于或等于~$n+1$。
\end{defn}

注意，一个公式或项可能有不止一个最小构造序列，但它们包含的符号串完全相同。

\begin{prop}
\ollabel{prop:subformula-equivs}
下列命题等价：
\begin{enumerate}
    \item $!B$ 是~$!A$ 的子公式。
    \item $!B$ 出现在~$!A$ 的每一个构造序列中。
    \item $!B$ 出现在~$!A$ 的某个最小构造序列中。
\end{enumerate}
\end{prop}

\begin{proof}
练习。
\end{proof}

\begin{prob}
证明 \olref[fol][syn][fseq]{prop:subformula-equivs}。
\end{prob}

\begin{history}
构造序列由 Raymond Smullyan（雷蒙德·斯穆里安）在其教科书 \emph{First-Order Logic} \citep{Smullyan1968} 中引入。
构造序列的更多性质由 \citet{Zuckerman1973}（祖克曼）建立。
\end{history}

\end{document}